% TÀI LIỆU / BÀI TẬP CẤP 3, dùng mẫu con mỏng toan10.sty (có danh mục lớp 10, nên dùng được \loai, \chuong).
\documentclass[12pt]{article}
\usepackage{toan10}
\begin{document}
\loai{1}\chuong{1}
\tieude{\tentailieu}{Bài 1. Mệnh đề}

\phan{A. Lý thuyết}
\begin{kienthuc}
Mệnh đề là một câu khẳng định đúng hoặc sai.
\end{kienthuc}

\phan{B. Phân loại và phương pháp giải toán}
\dang{Xét tính đúng sai của mệnh đề}
\muctl
\begin{bt}[TH]
Mệnh đề ``$2$ là số nguyên tố chẵn'' đúng hay sai?
\loigiai{Đúng.}
\end{bt}
\muctn
\begin{tn}[NB]
Câu nào là mệnh đề?
\pa{\dung{$2+3=5$.}}{Hôm nay trời đẹp quá!}{Bạn đi đâu đấy?}{$x+1>0$.}
\end{tn}
\mucds
\begin{ds}
Xét các mệnh đề về số nguyên tố.
\ydung{$7$ là số nguyên tố.}
\ysai{$9$ là số nguyên tố.}
\ydung{$2$ là số nguyên tố.}
\ysai{$1$ là số nguyên tố.}
\end{ds}
\muctln
\begin{tln}
Có bao nhiêu số nguyên tố nhỏ hơn $10$?
\traloi{4}
\end{tln}
\bangdapan
\end{document}
